Economists study how individuals and businesses make choices among alternative uses of scarce resources. This course introduces students to how markets align the interests of individuals and businesses to efficiently allocate the economy's scarce resources. Students will learn about supply, demand, and market equilibrium, as well as how economists study the welfare implications of government policies that alter market outcomes.

With the core tools of economic analysis at hand, students turn their attention to the subdiscipline of macroeconomics. Macroeconomists analyze the aggregate economic performance of nations. As macroeconomists do, the course separates the study of long-run economic growth from analysis of the short-run fluctuations of the business cycle. To engage in macroeconomics, students learn how macroeconomists measure key economic indicators like gross domestic product (GDP), inflation, and unemployment. Further, students learn about the key macroeconomic institutions in the U.S. and how to separately evaluate monetary and fiscal policy actions by applying the economic model of aggregate supply and aggregate demand.
